% Part: turing-machines
% Chapter: machines-computations
% Section: variants

\documentclass[../../../include/open-logic-section]{subfiles}

\begin{document}

\olfileid{tur}{mac}{var}
\olsection{اختلاف صيغ آلات تورنغ}

ليس تعريف آلة تورنغ مقصورًا على ما اخترناه؛ بل له صيغ كثيرة. فتعريفنا أوسع من غيره في أمور: نجيز أبجدية منتهية كيف اتفقت، وقد يقتصر غيرنا على $\TMstroke$ و~$\TMblank$؛ ونجيز الكتابة والحركة في خطوة واحدة، وقد يخص غيرنا الخطوة بأحدهما؛ ونجيز الكتابة مع سكون الرأس، وقد يحذف غيرنا «تعليمة» $\TMstay$. وهو أضيق من غيره في أمور أيضًا؛ إذ جعلنا الشريط ممتدًّا في جهة واحدة، وقد يجعل غيرنا له امتدادًا غير منته في الجهتين. بل يمكن اتخاذ عدة أشرطة، أو شبكة خانات لا نهائية. وجعلنا التعليمات دالة انتقال، وقد تجعل علاقة انتقال تبيح أكثر من تعليمة في الوضع الواحد.

والتوسيع الأخير ذو شأن خاص. ففي تعريفنا، إذا كانت الآلة في~$\OLId{latin-ordinary}{q}$ تقرأ~$\sigma$، عينت $\delta(\OLId{latin-ordinary}{q},\sigma)$ الرمز والحالة والموضع التاليين. أما إذا جعلنا التعليمات علاقة بين الزوج الراهن $\tuple{\OLId{latin-ordinary}{q},\sigma}$ والثلاثية الجديدة $\tuple{\OLId{latin-ordinary}{q}',\sigma',\OLId{latin-looped}{D}}$، فقد لا يتعين الفعل تعيينًا وحيدًا؛ إذ قد تضم العلاقة $\tuple{\OLId{latin-ordinary}{q},\sigma,\OLId{latin-ordinary}{q}',\sigma',\OLId{latin-looped}{D}}$ و$\tuple{\OLId{latin-ordinary}{q},\sigma,\OLId{latin-ordinary}{q}'',\sigma'',\OLId{latin-looped}{D}'}$ معًا. فهذه آلة تورنغ \emph{غير حتمية}، ولها مكان مهم في نظرية التعقيد الحسابي.

وكذلك اختلف الاصطلاح في التوقف: فعندنا تقف الآلة إذا لم تعرف دالة الانتقال، وعند غيرنا لا تقف إلا في حالة مخصوصة. وقد بينا في \olref[hal]{sec} أن تخصيص الحالة لا يغير ما تحسبه الآلة. ومن آثار امتداد شريطنا في جهة واحدة أنه قد تؤمر الآلة بما لا يتم على وجهه، وهو التحرك يسارًا وهي على أقصى خانة يسارًا. فجعلنا الرأس يبقى مكانه، ولا «يسقط» عن الشريط؛ وكان يمكن أن نجعل الآلة تقف. والتعريفان متكافئان: لمحاكاة آلة تقف عند محاولة التحرك يسارًا من~$\OLMathNumeral{0}$ نحذف كل انتقال $\delta(\OLId{latin-ordinary}{q},\TMendtape)=\tuple{\OLId{latin-ordinary}{q}',\sigma,\TMleft}$؛ فتقف بدل محاولة تجاوز~$\TMendtape$ يسارًا.\footnote{هذا الوصف ليس تامًّا \emph{على إطلاقه}؛ فقد يكتب $\TMendtape$ ويقرأ في غير الخانة~$\OLMathNumeral{0}$ (انظر \olref[una]{ex:mover}). ونتجاوز ذلك بتخصيص رمز ثان $\TMendtape'$ لهذه الاستعمالات.}

ويختلف أيضًا تمثيل الأعداد، فتختلف صورة دالة الدخل والخرج. أخذنا بالتمثيل الأحادي، ويجوز الأخذ بالثنائي؛ وحينئذ نحتاج إلى رمزين زائدين على $\TMblank$ و~$\TMendtape$.

غير أن هذه الفروق كلها لا تغير صنف الدوال القابلة للحساب بتورنغ. وهذه حقيقة لافتة: \emph{كل دالة تقبل الحساب بتورنغ في إحدى هذه الصيغ تقبله في سائرها}.

ولا نستوفي هنا براهين ذلك، بل نضرب مثالًا واحدًا: لا تزيد الأشرطة المتعددة، ولا الشريط الممتد في جهتين، القدرة الحسابية. ووجهه إجمالًا أن ذا الشريط الواحد الممتد في جهة واحدة يحاكي تلك الأشرطة. فنزيد حالات وتعليمات «تترجم» برنامج الآلة متعددة الأشرطة أو ذات الامتدادين إلى برنامج من صورتنا. فمثلًا نجعل الخانات الزوجية للشريط~$\OLMathNumeral{1}$، أو للجانب «الموجب» من الشريط ذي الجهتين، ونجعل الفردية للشريط~$\OLMathNumeral{2}$، أو للجانب «السالب».

\end{document}
